\section{Introduction}

High-quality training data remains the primary driver behind the empirical success of large language models (LLMs)~\citep{Hsieh2023,Muennighoff2023,Mohan2024,villalobos2024position}. 
Recent work has shown that intrinsic properties of training data, including data quality, distributional composition, and sample difficulty, shape the learning trajectory and downstream generalization performance of LLMs~\cite{Xu2023HardSample,Hoffmann2022Chinchilla,doremi}.
Consequently, a plethora of data-centric strategies, including selection, filtering, and reweighting, have been proposed to curate more effective training distributions~\citep{Sorscher2022Beyond,Xie2023DataSelection,Albalak2024Survey}.
Despite their disparate formulations, these methods share a common objective: to enhance model robustness and generalization by optimizing over a set of admissible data distributions $\mathcal{Q}$, rather than the raw empirical distribution~\cite{Sagawa2019GroupDRO}.

However, a key mechanism governing data utility, which involves the dynamic interactions of the train-unit (defined as inter-batch or inter-sample interactions, where the latter is a degenerate case of the former with a batch size of 1) interactions that evolve throughout training, has been largely overlooked. The efficacy of a train-unit depends not only on its intrinsic properties but also on when it is presented relative to others, as early train-units can prime the model for later ones, creating dependencies that render the data non-exchangeable. Recent empirical studies confirm that data ordering substantially influences convergence and final performance in LLM training~\cite{Kim2024StrategicOrdering,Pouransari2024DatasetDecomposition,Jia2025CurriculumEffects}. Yet these scheduling approaches remain largely heuristic; they lack a principled, optimization-grounded method to quantify the very interactions they seek to exploit. This leads to a core open question: how can we formally model dynamic train-unit interactions to derive optimal data schedules?

% From this perspective, a key limitation of prior work is the primary focus on refining global data distributions while overlooking the complex, intrinsic interactions between individual samples. 
% The existence of such interactions precludes the assumption of data exchangeability; instead, it dictates an inherently optimal training trajectory determined by the underlying dependencies between samples.
% Recent empirical studies challenge this assumption, demonstrating that data ordering can substantially influence both convergence behavior and final performance in LLM training~\cite{Kim2024StrategicOrdering,Pouransari2024DatasetDecomposition,Jia2025CurriculumEffects}.
% However, these approaches do not yet offer a principled understanding of why training order matters from an optimization dynamics perspective. 

To address this gap, we introduce $D^3$, a framework that reformulates data scheduling as a constrained optimization problem over a dynamic graph. We start from the principle that the optimal order should prioritize train-units that most effectively “prepare” the model for subsequent data. $D^3$ makes this precise by modeling pairwise influence via look-ahead loss reduction: if training sample $i$ now significantly lowers the future loss of train-unit $j$, then $i$ should precede $j$. By estimating these directional effects throughout training, $D^3$ maintains a dynamic influence graph that captures the evolving topology of train-unit dependencies. The training sequence is then derived as the solution that best respects this graph’s structure, ensuring the model follows a coherent, interaction-aware learning trajectory. Our contributions are summarized as follows:  
\begin{itemize}
    \item We propose $D^3$, a data scheduling framework that formalizes sequence ordering as an optimization problem over a dynamic training influence graph.
    \item By providing a fundamental method to characterize asymmetric train-unit interactions, we reveal that these relationships dictate the necessity of a structured training sequence, addressing a significant gap in conventional distribution-centric approaches.
    \item Extensive numerical validations across multiple pre-training and post-training tasks show that $D^3$ consistently improves the performance of LLMs with reasonable computation overheads. 
\end{itemize}